\documentclass[11pt]{article}
\usepackage[a4paper,margin=2.15cm]{geometry}
\usepackage{lmodern}
\usepackage{amsmath,amssymb,amsthm,mathtools}
\usepackage{booktabs}
\usepackage{microtype}
\usepackage[hidelinks]{hyperref}
\usepackage{graphicx}
\graphicspath{{figs/}}
\usepackage{enumitem}
\usepackage{titlesec}
\titlespacing*{\section}{0pt}{1.3ex plus .2ex}{0.7ex}
\titlespacing*{\subsection}{0pt}{1.1ex plus .2ex}{0.5ex}

\theoremstyle{plain}
\newtheorem{theorem}{Theorem}[section]
\newtheorem{lemma}[theorem]{Lemma}
\newtheorem{proposition}[theorem]{Proposition}
\theoremstyle{definition}
\newtheorem{definition}[theorem]{Definition}
\newtheorem{remark}[theorem]{Remark}
\newtheorem{example}[theorem]{Example}

\newcommand{\R}{\mathbb{R}}
\newcommand{\eps}{\varepsilon}
\newcommand{\sgn}{\operatorname{sgn}}
\newcommand{\htop}{h_{\mathrm{top}}}
\setlist{itemsep=2pt,topsep=3pt,leftmargin=1.5em}
\setlength{\abovedisplayskip}{5pt}
\setlength{\belowdisplayskip}{5pt}
\setlength{\abovedisplayshortskip}{3pt}
\setlength{\belowdisplayshortskip}{4pt}


\begin{document}

\begin{center}
  {\Large\bfseries Bands and renormalization in the logistic family}\\[3pt]
  {\small companion to \texttt{core\_entropy\_notes.tex}}
\end{center}
\vspace{2pt}

\section{Setting}\label{sec:setting}

These notes are a companion to \texttt{core\_entropy\_notes.tex}, which computes the core
entropy $h(\lambda)$ of the logistic map $g_\lambda(x)=\lambda x(1-x)$, $x\in[0,1]$,
$\lambda\in[2,4]$, as the reciprocal of the smallest zero of the kneading determinant.  The
derivation, the algorithm and its error control are there; here we use that machinery and
that notation.  Let $\nu=\nu_0\nu_1\cdots$ be the itinerary of the critical point, each term
$0$ or $1$ according to the side of $\tfrac12$ on which $g_\lambda^k(\tfrac12)$ falls, and
the return symbol $*$ when it falls exactly on $\tfrac12$; let $\eps_k=\pm1$ be the signs of
$g_\lambda'$ along the critical orbit, a return to $\tfrac12$ being read as the symbol $1$;
and let
\[
  D(t)=\sum_{k\ge0}\eps_k t^k
\]
be the kneading determinant.  Its smallest zero in $(0,1)$ is $t_0$, the entropy is
$h(\lambda)=\log(1/t_0)$, and $s=e^{h}=1/t_0$ is the slope of the piecewise-linear model
$f_s$, the tent map to which $g_\lambda$ is semi-conjugate.  We write $c_n$ for the number of
solutions of $g_\lambda^n(x)=\tfrac12$, which grows like $s^n$; it is the modulation of
$c_ns^{-n}$, rather than the count itself, that carries the information about the bands.  Let
$\lambda_\infty=3.5699456\ldots$ be the accumulation point of the period-doubling cascade.  A
parameter is a \emph{Misiurewicz parameter} when its critical point is strictly preperiodic,
that is, when $g_\lambda^k(\tfrac12)=g_\lambda^j(\tfrac12)$ for some $k>j\ge0$.

What follows concerns the \emph{limit set} of $g_\lambda$ --- the set of limit points of a
typical orbit, that is, the attractor of the map --- and in particular what its components
look like when it is not an interval.  Three examples are worked out exactly, and the last
section draws the corresponding complex picture.

\section{Bands and renormalization}\label{sec:bands}

The limit set of $g_\lambda$ is not always an interval, and the pattern of its components is
visible in the determinant.  For typical $\lambda>\lambda_\infty$ the limit set of a typical
orbit is a union of $2^k$ closed intervals, the \emph{bands}; as $\lambda$ increases the bands
merge pairwise, and at $\lambda=3.6785735\ldots$ --- the Misiurewicz parameter, where the critical orbit
lands on the fixed point $1-1/\lambda$ --- the last two merge and the limit set becomes a single interval.  The mergers are the parameters at
which $s=e^{h(\lambda)}$ takes the values $2^{1/2^k}$, that is $h(\lambda)=2^{-k}\log2$, and
they are exactly the thresholds at which $c_ns^{-n}$ stops oscillating with period $2^k$,
as the counts of the companion note show.  In the determinant they appear as the modulus of its zeros: for
$\lambda$ in the $2^k$-band range $D$ has $2^k$ zeros in $|t|<1$, all of modulus $t_0$, at
$t_0\zeta$ for the $2^k$-th roots of unity $\zeta$.

The thresholds $s^{2^k}=2$ are the band-merging parameters, and the exact example below shows
in miniature why the determinant sees them: in the four-band case the numerator of $D$ is a
polynomial in $t^4$, so its zeros sit at $t_0$ times the fourth roots of unity, and each band
carries the same period-$3$ dynamics.

The two counts of $2^k$ are one, in the period-doubling cascade.  There the bands are the intervals of
the $2^k$-fold doubling renormalization, the map on a band is $f_s^{2^k}$ with slopes
$\pm s^{2^k}$ and is a full tent map exactly at the merger $s^{2^k}=2$, and the band
determinant has the single zero $t_0^{2^k}$.  So the $2^k$ zeros of $D$ at $t_0\zeta$ are the
$2^k$ determinations of that one zero, since $t\mapsto t^{2^k}$ takes $\zeta$ to $1$;
equivalently $c_ns^{-n}$ is a sum of $2^k$ interleaved sequences, one for each residue class
of $n$ modulo $2^k$, and converges only along those classes --- the period-$2^k$ modulation is
the band structure read off the counts.  This is Corollary~9.5 of \cite{MT}: the lap
numbers differ from a finite sum $\sum_j a_j^{-k}$ over the zeros of $D$ in the disc by
$O\bigl((1+\varepsilon)^k\bigr)$, so a set of zeros on the circle $|t|=t_0$ leaves a main term
periodic in $k$; that is the ``pure periodicity'' of \cite{FS}, \S IV.6.1, where dominant singularities
regularly spaced on the circle of convergence produce a superposition of regimes indexed by the
residue class of the index.  They give the case $k=1$ --- zeros at $\pm1/s$, and lap numbers with one
asymptotic for $k$ even and another for $k$ odd --- and their Example~14.4 is the
substitution $D\mapsto(1-t)D(t^2)$ that produces the zeros.  Outside the cascade the
bookkeeping changes: in the period-$3$ copy of \S\ref{sec:bands3} the band zeros are
subdominant, though the factor $t\mapsto t^p$ that produces $p$ bands is the same.

\subsection{The periodicity of the counts}\label{sec:periodicity}

Plotted as a single sequence, $\log c_n-h(\lambda)n$ zigzags; split by the residue of $n$
modulo $2^k$ it falls into $2^k$ smooth curves, one for each band.  Figure~\ref{fig:lapsperiod}
shows this at two parameters, one in the two-band range and one in the four-band range.

\begin{figure}[htbp]
  \centering
  \includegraphics[width=\textwidth]{fig_laps_period}
  \caption{The periodicity of the counts.  $\log c_n-h(\lambda)n$ against $n$, with the points
  coloured by the residue of $n$ modulo $2^k$: (a)~the period-$6$ centre $\lambda=3.6276$, in
  the two-band range, gives two curves; (b)~the period-$12$ centre $\lambda=3.5822$, the
  parameter of \S\ref{sec:bands4}, in the four-band range, gives four.  Each colour is smooth
  while the union zigzags.}
  \label{fig:lapsperiod}
\end{figure}

\subsection{An exact example with four bands}\label{sec:bands4}

At $\lambda_c=3.5822298358\ldots$ the critical point has period $12$; $\lambda_c$ is the
centre of a period-$12$ component inside the four-band window $(3.5748,3.5926)$.  Writing
$x_j=g_{\lambda_c}^j(\tfrac12)$, the orbit is
\[
  \begin{array}{llll}
  x_0=0.5 & x_1=0.895557459 & x_2=0.335061348 & x_3=0.798103760\\
  x_4=0.577219553 & x_5=0.874197126 & x_6=0.393961138 & x_7=0.855278006\\
  x_8=0.443399590 & x_9=0.884081405 & x_{10}=0.367112196 & x_{11}=0.832298259
  \end{array}
\]
and $\nu=(1,0,1,1,1,0,1,0,1,0,1,1)$ has period $12$, as does $\eps$.  The determinant collapses
to
\begin{equation}\label{eq:twelve}
  D(t)=\frac{(t-1)(t^8+t^4-1)}{(1+t^2)(1+t^4+t^8)},
\end{equation}
whose poles all lie on $|t|=1$ and whose zeros in the disc are the four roots of
$t^4=\tfrac1\tau$:
\[
  t_0=\tau^{-1/4}=0.8866517793\ldots,\qquad h(\lambda_c)=\tfrac14\log\tau=0.1203029563\ldots
\]
The other four roots of $t^8+t^4-1$ have modulus $\tau^{1/4}>1$ and lie outside.  The
substitution $u=t^4$ turns the numerator into $u^2+u-1$, the polynomial of the period-$3$ example of the companion note: each band carries a period-$3$ renormalization with entropy
$\log\tau$, and the entropy of the whole map is a quarter of it.

The bands are visible at a nearby parameter.  At $\lambda_c-10^{-3}$ the attractor consists of
the four intervals
\[
  [0.3357,0.3927],\quad[0.4463,0.5760],\quad[0.7986,0.8541],\quad[0.8746,0.8953],
\]
whose eight endpoints are the first eight images of the critical value, so they are exact rather
than measured.  At $\lambda_c$ the twelve points fall three at a time into their limits, by
residue modulo four:
\[
  B_r=\{x_j:\ j\equiv r \ (\mathrm{mod}\ 4)\},
  \qquad r=0,1,2,3 .
\]
Then $g_\lambda$ maps $B_r$ into $B_{r+1}$, cyclically, and $g_\lambda^4$ maps each band into
itself as a unimodal map whose critical orbit is the $3$-cycle in that band --- the
renormalized map, with its own period-$3$ centre.  Figure~\ref{fig:bands} shows the two
pictures together.

\begin{figure}[htbp]
  \centering
  \includegraphics[width=\textwidth]{fig_bands}
  \caption{The four bands.  The shaded density is the attractor at $\lambda_c-10^{-3}$, a
  union of four intervals; the coloured points are the twelve points of the critical orbit at
  $\lambda_c$, grouped by band as $x_j$ with $j\equiv r \pmod 4$.  Each group forms a
  $3$-cycle of $g_\lambda^4$ on its band.}
  \label{fig:bands}
\end{figure}

\subsection{A three-band example}\label{sec:bands3}

The period-$3$ copy behaves differently, and a three-band example makes the difference clear.
Its chaotic part lies between the accumulation of its own doublings, at
$\lambda\approx3.8495$, and the crisis at $\lambda\approx3.8568$ where the three bands merge
into one; beyond $3.8568$ the attractor is a single interval, so the bands live between those
two values.

At $\lambda_9=3.8536752768\ldots$ the critical point has period $9=3\cdot3$.  With
$x_j=g_{\lambda_9}^j(\tfrac12)$,
\[
  \begin{array}{lll}
  x_0=0.5 & x_1=0.963418819 & x_2=0.135815070\\
  x_3=0.452303312 & x_4=0.954651808 & x_5=0.166832284\\
  x_6=0.535658061 & x_7=0.958518881 & x_8=0.153223807
  \end{array}
\]
and $\nu=(1,0,0,1,0,1,1,0,*)$ has period $9$.  The determinant is
\begin{equation}\label{eq:nine}
  D(t)=\frac{(t^2+t-1)(t^6+t^3-1)}{(1+t)(t^2-t+1)(t^6-t^3+1)}
      =\frac{t^8+t^7-t^6+t^5+t^4-t^3-t^2-t+1}{1+t^9},
\end{equation}
whose zeros in the disc are $t_0=1/\tau=0.618\ldots$, from the factor $t^2+t-1$, and the three
roots of $t^3=\tfrac1\tau$ at $0.8518\ldots$ times the cube roots of unity, from
$t^6+t^3-1=(t^3)^2+t^3-1$.  So $h(\lambda_9)=\log\tau$, the value the whole period-$3$ copy
carries, and the band structure sits in the second factor: the substitution $t\mapsto t^3$
turns it into the same golden polynomial, exactly as $t\mapsto t^4$ did at the period-$12$
centre.

The nine points fall three at a time into the bands $B_r=\{x_j:\ j\equiv r\pmod3\}$, whose
hulls are
\[
  B_0=[0.4523,0.5357],\qquad B_1=[0.9547,0.9634],\qquad B_2=[0.1358,0.1668],
\]
with $\tfrac12\in B_0$; $g_\lambda$ maps $B_r$ into $B_{r+1}$ cyclically, and $g_\lambda^3$
carries each band into itself with a superattracting $3$-cycle.

One difference from the four-band case is worth recording.  There $t_b^{1/4}=t_0$ was the
smallest zero of $D$, so the zeros at the fourth roots of unity sat on the circle of
convergence and the period-$4$ modulation of the counts did not decay.  Here $t_b^{1/3}=0.8518$
is larger than $t_0=0.618$, so the three cube-root zeros are subdominant and the period-$3$
modulation decays like $(t_0/t_b^{1/3})^n=\tau^{-2n/3}$; the counts settle onto a line after
all, only more slowly.

\subsection{The merging parameter}\label{sec:bandsmerge}

At a Misiurewicz value the two pictures degenerate into one, and the determinant becomes as
simple as it can be.  At $\lambda_M=3.6785735104\ldots$ the critical orbit is
\[
  \tfrac12\longmapsto0.9196433776\longmapsto0.2718445063\longmapsto q,
  \qquad q=1-\tfrac1{\lambda_M}=0.7281554937,
\]
with $q$ a repelling fixed point, so $\nu=(1,0,1,1,1,\ldots)$ and
\begin{equation}\label{eq:merging}
  D(t)=\frac{1-2t^2}{1+t},
\end{equation}
whose zeros are $t=\pm1/\sqrt2$: two of them, of modulus
$t_0=1/\sqrt2$, as the two bands require, with $h=\log\sqrt2=0.3465736\ldots$  The gap between
the bands has closed at $q$ --- at $\lambda_M-2\cdot10^{-3}$ the attractor still avoids a
neighbourhood of $q$, at $\lambda_M$ it does not --- but the merging survives in the invariant
density, which has a square-root singularity at $q$: measured by a histogram, the excess over
the background in a bin of width $w$ centred at $q$ grows by $2.4$ when $w$ is reduced by $5$,
as $|x-q|^{-1/2}$ predicts ($\sqrt5=2.24$).  So the two bands are still visible, as the
lifetime of an orbit in a neighbourhood of the point where they meet.

\section{The complex picture}\label{sec:julia}

The conjugacy of the companion note, $z=\psi_\lambda(x)=\lambda(\tfrac12-x)$, carries
$g_\lambda$ to $z\mapsto z^2+c$ with $c=\lambda/2-\lambda^2/4$, and carries $[0,1]$ onto
$[-\lambda/2,\lambda/2]$.  So every statement about $g_\lambda$ is one about a quadratic
polynomial, and it is natural to ask what the bands look like there.  Let the \emph{filled
Julia set} of $z^2+c$ be the set of $z$ whose orbit stays bounded.  On the real axis the
filled set is exactly the segment $[-\lambda/2,\lambda/2]$: the endpoint $\lambda/2$ is the
fixed point $\beta$ of $z^2+c$ and $-\lambda/2$ is its preimage, since $c$ and $\lambda$
satisfy $\beta^2+c=\beta$ with $\beta=\lambda/2$.  Off the real axis the structure that Figure~\ref{fig:julia}
draws is the Julia set, the boundary of the filled set: not an interval, but a continuum whose
symmetrically placed arms hang from that spine.  At the two centres the filled set has
interior --- the Fatou components of the superattracting cycles --- so the two sets differ
there; at the merging parameter the critical point is preperiodic, the map is expanding on its
Julia set, and the filled set is the Julia set itself.

\begin{figure}[htbp]
  \centering
  \includegraphics[width=\textwidth]{fig_julia}
  \caption{The filled Julia set of $z^2+c$ at the three parameters of
  \S\ref{sec:bands4}, \S\ref{sec:bands3} and \S\ref{sec:bandsmerge}, with
  $c=\lambda/2-\lambda^2/4$: (a)~the period-$12$ centre, $c=-1.41698$; (b)~the period-$9$
  centre, $c=-1.78587$; (c)~the merging parameter, $c=-1.54369$.  The colour encodes the
  escape time of $z$ on a logarithmic scale, cycled nine times, so that the level sets appear
  as coloured rings; black is the filled set, where the orbit never leaves.  The window is
  $[-2.05,2.05]^2$ throughout, so the panels are drawn to the same scale, and the middle row
  of each panel is the real axis: it meets the filled set in $[-\lambda/2,\lambda/2]$, that is
  in $[-1.7911,1.7911]$, $[-1.9268,1.9268]$ and $[-1.8393,1.8393]$.  The yellow marks on that
  row are the limit set of a typical orbit, carried to the plane by the conjugacy: the twelve
  points of the attracting cycle at (a), the nine at (b), and at (c) the single interval
  $[c,c^2+c]=[-1.5437,0.8393]$ into which the last two bands have merged.}
  \label{fig:julia}
\end{figure}

One thing to know before reading the figure is that the Julia set has no area: writing
$\mathcal{H}^2$ for two-dimensional Hausdorff measure and $\dim_H$ for Hausdorff dimension, $\dim_HJ<2$ for
all three of these maps --- the two centres because they are hyperbolic, the merging parameter
because it is expanding on its Julia set --- and hence $\mathcal{H}^2(J)=0$, so the filament
has measure zero however far it is resolved.  The filled set does have area, but only at the
two centres: there the Fatou components make up $3\cdot10^{-4}$ and $6\cdot10^{-5}$ of the
window, so at the scale of the figure they are no more than specks.  So the black is not where the structure is --- at
(a)~and~(b) it is those components, at~(c) there is next to none of it --- and what the eye
follows is the Julia set, drawn by the escape-time bands.

The yellow marks then answer the question the panel titles raise, and expose a discrepancy in
them.  At the two centres the limit set is the attracting cycle itself --- twelve
points at~(a), nine at~(b) --- so it is a finite set there, and the bands of those titles are
not present at the parameters drawn: they are what the cycle becomes on moving off the
centre, the four intervals of \S\ref{sec:bands4} at $\lambda_c-10^{-3}$.  At~(c) the limit set
is the interval $[c,c^2+c]=[-1.5437,0.8393]$, of length $2.383$, and that is what the bar
shows.  In all three the limit set is one-dimensional, and the complex picture around it is
not a second view of the bands but the geometry in which their renormalization of
\S\ref{sec:bands4} and \S\ref{sec:bands3} lives.

\begin{thebibliography}{9}\footnotesize\setlength{\itemsep}{1pt}\setlength{\parskip}{0pt}
\bibitem{MT} J.~Milnor and W.~Thurston, On iterated maps of the interval, in
\emph{Dynamical Systems} (College Park, 1986--87), Lecture Notes in Math. 1342,
Springer, 1988, 465--563.
\bibitem{FS} P.~Flajolet and R.~Sedgewick, \emph{Analytic Combinatorics},
Cambridge University Press, 2009, \S IV.6.
\end{thebibliography}

\end{document}
